
\documentclass[12pt]{article}
\usepackage{amsmath}
\usepackage{geometry}
\geometry{a4paper, margin=1in}

\title{Transcript for Presentation on Difference Equations}
\author{Your Name}
\date{\today}

\begin{document}
\maketitle

\section*{Slide 1: Overview of Difference Equations}
Welcome! In this presentation, we’ll be exploring the concept of \textbf{difference equations}. 

Difference equations are widely used for \textbf{modeling dynamic systems} by tracking the relationship between values at different time steps.

Next, let’s discuss some \textbf{key examples and applications}. These include \textbf{control systems}, \textbf{signal processing}, and \textbf{finance}.

Each of these fields relies on the precise modeling of discrete time steps.

Finally, difference equations share a close \textbf{relationship with differential equations}, which describe continuous-time systems. In many cases, difference equations provide a discrete approximation.

\section*{Slide 1a: Importance of Difference Equations}
Why are difference equations so important?

First, they’re essential for \textbf{modeling dynamic systems}. They capture changes between discrete time steps—ideal for systems like population growth or economic trends.

Next, they serve as a \textbf{bridge to differential equations}. When data is sampled at discrete intervals—such as digital signals—difference equations offer a way to model those systems effectively.

\section*{Slide 1b: Applications of Difference Equations}
Here are some key \textbf{application areas} of difference equations.

In \textbf{control systems}, they’re used in feedback loops, system stability, and controller design.

In \textbf{signal processing}, difference equations help with filtering and analyzing discrete-time signals.

In \textbf{finance}, they model investments, loans, and asset dynamics over time.

And finally, in \textbf{epidemiology}, they predict the spread of diseases using discrete time models.

\section*{Slide 2: Definitions (Scalar and Vector Case)}
Let’s look at the \textbf{definition of difference equations}.

A \textbf{scalar difference equation} relates successive values of a scalar sequence. Here’s an example:
\[
y(n+1) = a(1) y(n) + a(2) y(n-1) + \cdots + a(k) y(n-k+1) + b.
\]
When the output becomes a \textbf{vector}, we use a \textbf{vector difference equation}:
\[
y_i(n+1) = \sum_{j=1}^{k} a_{ij} y_j(n) + b_i.
\]
Scalar equations are a special case of vector equations when the output has only one dimension.

\section*{Slide 2b: Higher-Order vs. First-Order Difference Equations}
Higher-order scalar equations involve \textbf{multiple past values} of the sequence, such as \( y(n) \) and \( y(n-1) \).

However, these can be converted into \textbf{first-order vector equations} involving only the current and next time steps. For example:
\[
x(n+1) = A x(n) + B u(n).
\]

\section*{Slide 2c and 2d: Converting Higher-Order Scalar to First-Order Vector Equations}
Let’s see how to convert higher-order scalar equations into first-order vector equations.

Start with the scalar equation:
\[
y(n+1) = a(1) y(n) + a(2) y(n-1) + b.
\]
We introduce \textbf{state variables} to capture past steps:
\[
x_1(n) = y(n), \quad x_2(n) = y(n-1).
\]
Now, we can express the system using \textbf{state-space equations}:
\[
x_1(n+1) = a(1) x_1(n) + a(2) x_2(n) + b, \quad x_2(n+1) = x_1(n).
\]

\section*{Slide 3a: State Variables and Outputs}
In state-space models, the \textbf{state vector} \( x(n) \) represents the internal state at time \( n \). Its dimension is \( d \), the number of states.

The \textbf{output vector} \( y(n) \) provides the observed values. Its dimension is \( m \), the number of outputs.

\section*{Slide 3b: State-Space Output Equation}
The output \( y(n) \) is related to the state through the matrix \( C \):
\[
y(n) = C x(n).
\]
Here, \( C \) is an \( m \times d \) matrix, where \( m \) is the number of outputs and \( d \) is the number of states.

\section*{Slide 3c: Control Inputs, Noise, and State Transition}
Control inputs \( u(n) \) represent external influences on the system. These inputs have dimension \( p \).

Noise inputs \( w(n) \) account for uncertainties in the system, with dimension \( q \).

The state transition equation is:
\[
x(n+1) = A x(n) + B u(n) + G w(n).
\]

\section*{Slide 4a: Defining the State for a Simple System}
Let’s consider a simple system described by:
\[
y(n+1) = a y(n) + b u(n).
\]
We define the \textbf{state variable} \( x(n) = y(n) \). This makes the state and output identical, though in more general cases, they may differ.

\section*{Slide 4b: State-Space Representation (General Form)}
The general state-space form is:
\[
x(n+1) = A x(n) + B u(n).
\]
The output is:
\[
y(n) = C x(n).
\]
Here, \( A \), \( B \), and \( C \) are matrices of appropriate dimensions.

\section*{Slide 5a1: AR Component of ARMA Models}
The AR component of an ARMA model is:
\[
y(n) = \mu_{AR} + \sum_{i=1}^{p} a(i) y(n-i) + w(n).
\]
Here, \( \mu_{AR} \) is a constant drift term, \( a(i) \) are the AR coefficients, and \( w(n) \) is white noise.

\section*{Slide 5a2: MA Component and Combined ARMA Model}
The MA component is:
\[
y(n) = \mu_{MA} + w(n) + \sum_{j=1}^{q} b(j) w(n-j).
\]
The full ARMA model is:
\[
y(n) = \mu + \sum_{i=1}^{p} a(i) y(n-i) + w(n) + \sum_{j=1}^{q} b(j) w(n-j).
\]

\section*{Slide 5b and 5c: Explanation of Mu in ARMA Models}
In the AR component, \( \mu_{AR} \) shifts the series level.

In the MA component, \( \mu_{MA} \) adjusts the long-term mean. These terms combine into \( \mu = \mu_{AR} + \mu_{MA} \).

\end{document}
